\chapter{Evaluating a strategy}\label{ch:eval}

\noindent A backtest ends with a number; an evaluation starts with a battery. Sharpe is
necessary and nowhere near sufficient: it has a huge standard error at realistic sample
sizes, it ignores the path (drawdown), the shape (skew), and the bill (turnover). This
chapter is the suite.

\section{The Sharpe CI problem}\label{sec:sharpeci}

The standard error of a measured Sharpe ratio is roughly
$\mathrm{se}\approx\sqrt{(1+\mathrm{SR}^2/2)/T}$. \texttt{sharpe\_se.py}:
\begin{lstlisting}[numbers=none]
T=  63 days  se(SR)~0.154  95% CI on SR=1: [0.70,1.30]
T= 252 days  se(SR)~0.077  95% CI on SR=1: [0.85,1.15]
T= 756 days  se(SR)~0.045  95% CI on SR=1: [0.91,1.09]
T=2520 days  se(SR)~0.024  95% CI on SR=1: [0.95,1.05]
empirical se over 2000 sims (T=252): 1.012
\end{lstlisting}

\begin{example}\label{ex:sharpeci}
\textbf{The one-year Sharpe is a rumor.} A year of daily data narrows a true SR-1
strategy only to $[0.85,1.15]$ in theory---and the empirical spread across 2{,}000
sims is $\pm1.0$ (the non-normal world is wider still).
\begin{solution}
At $T{=}252$ you cannot distinguish Sharpe 1.5 from Sharpe 0.3 with any confidence.
The discipline consequence: small Sharpe differences between two strategies on one
year of data are noise, and a single-year backtest ``proving'' Sharpe 1.8 is compatible
with true Sharpe anywhere in $[0.5,3]$---before the multiple-testing haircut from ch09.
\end{solution}
\end{example}

\begin{figure}[htbp]\centering
\begin{tikzpicture}[font=\scriptsize]
  \draw[->,slate] (0,0)--(0,4.4);
  \foreach \y/\t/\lo/\hi in {3.8/63d/0.70/1.30, 2.6/252d/0.85/1.15, 1.4/756d/0.91/1.09, 0.4/2520d/0.95/1.05}{
    \node[left] at (-0.1,\y) {\t};
    \draw[accent,thick] (\lo*3,\y)--(\hi*3,\y);
    \fill[ink] (3,\y) circle (1.6pt);}
  \draw[slate,dashed] (3,-0.3)--(3,4.6) node[above]{SR $=1$};
  \node[below] at (3,-0.4) {95\% CI on measured SR (true SR $=1$, normal)};
\end{tikzpicture}

\caption{Sharpe confidence intervals shrink only as $1/\sqrt T$---a decade buys a
$\pm5\%$ band on the truth.}
\label{fig:ci}
\end{figure}

\section{The metric suite}\label{sec:suite}

\texttt{metrics\_demo.py} runs two strategies with nearly identical hit rates:
\begin{lstlisting}[numbers=none]
tail-risk strategy: sharpe +1.49  mdd -7.9%  hit 58.4%
boring strategy:    sharpe +3.63  mdd -4.6%  hit 59.7%
same mean, different product -- the suite shows what sharpe hides
\end{lstlisting}

\begin{example}\label{ex:suite}
\textbf{Why the suite, not the number.} The boring strategy's Sharpe is 2.4$\times$ the
tail-risk one's---the hit rates are within a point of each other.
\begin{solution}
Hit rate alone describes neither. The suite separates what one number merges:
Sharpe (mean/vol), MDD (worst path), skew (asymmetry), hit rate (frequency). A strategy
reporting one of them is choosing which weakness to hide.
\end{solution}
\end{example}

\begin{figure}[htbp]\centering
\begin{tikzpicture}
\begin{axis}[width=0.86\linewidth,height=0.32\linewidth,
    xlabel={day},ylabel={cum.\ log return},
    xmin=0,xmax=1160,
    tick style={draw=none},axis line style={draw=slate},font=\scriptsize]
  \addplot[accent,thick] file{figures/data/EQ12.tex};
\end{axis}
\end{tikzpicture}
\\[4pt]
\begin{tikzpicture}
\begin{axis}[width=0.86\linewidth,height=0.22\linewidth,
    xlabel={day},ylabel={dd},
    xmin=0,xmax=1160,ymin=-0.10,ymax=0,
    tick style={draw=none},axis line style={draw=slate},font=\scriptsize]
  \addplot[red,thick] file{figures/data/DD12.tex};
\end{axis}
\end{tikzpicture}

\caption{Equity curve and its drawdown series: the middle panel is where the risk lives.}
\label{fig:eqdd}
\end{figure}

\begin{figure}[htbp]\centering
\begin{tikzpicture}[font=\scriptsize,every node/.style={align=center}]
  \node[draw=ink,fill=goodbg,rounded corners=1pt,minimum width=2.5cm] (s) at (0,2.0) {Sharpe\\ (mean/sd)};
  \node[draw=slate,fill=white,rounded corners=1pt,minimum width=2.5cm] (m) at (3.2,2.0) {MDD\\ (path risk)};
  \node[draw=slate,fill=white,rounded corners=1pt,minimum width=2.5cm] (k) at (6.4,2.0) {skew / worst\\ (tail shape)};
  \node[draw=slate,fill=white,rounded corners=1pt,minimum width=2.5cm] (h) at (0,0.6) {hit rate\\ (payoff mix)};
  \node[draw=slate,fill=white,rounded corners=1pt,minimum width=2.5cm] (t) at (3.2,0.6) {turnover\\ (cost load)};
  \node[draw=slate,fill=white,rounded corners=1pt,minimum width=2.5cm] (c) at (6.4,0.6) {capacity\\ (ADV footprint)};
  \node[text=slate] at (3.2,-0.5) {no single number survives contact with the others};
\end{tikzpicture}

\caption{The evaluation suite: each metric answers a question the others cannot.}
\label{fig:metrics}
\end{figure}

\section{The cliff hidden in high Sharpe}\label{sec:cliff}

\texttt{tail\_demo.py} builds the classic put-seller profile---steady wins, rare crashes:
\begin{lstlisting}[numbers=none]
put-seller profile: sharpe +2.09  skew -4.30  worst -11.2%
high sharpe can hide a cliff: read skew and worst day too
\end{lstlisting}

\begin{example}\label{ex:cliff}
\textbf{Sharpe 2.09 with skew $-4.3$.} Ninety-five percent of days pay +0.7\%; 5\% lose
$-8\%$.
\begin{solution}
Volatility---Sharpe's denominator---is symmetric, so it treats the cliff the same as the
ladder. Skew ($-4.3$) and worst-day ($-11.2\%$) report what Sharpe cannot: the distribution
is a staircase down with a trapdoor. Any metric suite that excludes the left tail is an
invitation to sell insurance and call it alpha.
\end{solution}
\end{example}

\begin{figure}[htbp]\centering
\begin{tikzpicture}
\begin{axis}[width=0.86\linewidth,height=0.34\linewidth,
    xlabel={daily return},ylabel={count},
    tick style={draw=none},axis line style={draw=slate},font=\scriptsize]
  \addplot[fill=accent!30,draw=accent,thick] file{figures/data/TAILH.tex};
  \addplot[red,dashed,thick] coordinates {(-0.08,0)(-0.08,60)};
  \node[font=\scriptsize,text=red,align=left] at (axis cs:-0.16,55) {the cliff\\ (skew $-4.3$)};
\end{axis}
\end{tikzpicture}

\caption{The return histogram the Sharpe ratio summarizes as one number.}
\label{fig:tailhist}
\end{figure}

\begin{figure}[htbp]\centering
\begin{tikzpicture}[font=\scriptsize]
  \draw[->,slate] (0,0)--(8.6,0) node[below left=-1pt]{day};
  \foreach \x in {0.3,0.9,1.5,2.1,2.7,3.3,3.9,4.5,5.1,5.7,6.3,6.9,7.5}
    \draw[good,thick] (\x,1.0)--(\x,1.35);
  \draw[red,very thick] (6.9,1.0)--(6.9,-0.9);
  \node[align=left,text=slate] at (8.2,2.0) {+0.7\% on 95\% of days,\\ $-8\%$ on 5\%:\\ Sharpe $+2.09$};
\end{tikzpicture}

\caption{Steady small wins punctuated by a rare large loss---the profile that maximizes
Sharpe-per-unit-of-honesty.}
\label{fig:cliff}
\end{figure}

\section{Attribution}\label{sec:attr}

The last question: \emph{where} did the return come from? Signal (the rule's direction),
timing (when it traded), sizing (how much), costs (what it paid). A strategy whose PnL is
90\% one lucky sizing decision is not the strategy described.

\begin{example}\label{ex:attr}
\textbf{Attribution before celebration.} Suppose the backtest returns +22\%: signal
+35\%, sizing $-8\%$, costs $-5\%$.
\begin{solution}
The marketable product is the signal---the sizing loss is a fixable implementation
detail and the cost is a capacity ceiling. Attribution converts one number into a list
of actions: fix sizing, shrink footprint, keep the signal.
\end{solution}
\end{example}

\begin{figure}[htbp]\centering
\begin{tikzpicture}[font=\scriptsize]
  \draw[->,slate] (0,0)--(8.4,0) node[below left=-1pt]{return};
  \draw[fill=accent!30,draw=accent] (0.3,0) rectangle (2.3,1.1) node[midway]{signal};
  \draw[fill=amber!50,draw=amber] (2.3,0) rectangle (3.5,0.7) node[midway]{timing};
  \draw[fill=good!40,draw=good] (3.5,0) rectangle (4.3,0.45) node[midway]{sizing};
  \draw[fill=red!30,draw=red] (4.3,0) rectangle (5.1,0.25) node[midway]{costs};
  \draw[fill=slate!25,draw=slate] (5.1,0) rectangle (5.9,0.1) node[midway]{net};
  \node[align=left,text=slate] at (8.0,0.9) {attribute\\ before you\\ celebrate};
\end{tikzpicture}

\caption{Return decomposition: gross signal minus timing, sizing and costs is the net
the strategy actually ships.}
\label{fig:attr}
\end{figure}

\begin{checkpoint}[Evaluation audit]
\begin{itemize}
\item What is the 95\% CI on a measured Sharpe of 1.0 over one year? ($\approx$ two
Sharpe points wide---theory $\pm0.3$, empirically wider.)
\item Which two metrics expose a put-selling profile? (Skew and worst-day.)
\item Name the attribution categories. (Signal, timing, sizing, costs.)
\end{itemize}
\end{checkpoint}

\section{Recap}\label{sec:evalrecap}

Sharpe is a coarse instrument: wide confidence intervals, blind to path and tail. The
evaluation suite---MDD, skew, hit rate, turnover, capacity, attribution---is what
converts a backtest number into a decision. Report all of them or name the one you
are hiding.

\begin{table}[htbp]\centering\small
\caption{Chapter toolbox.}
\begin{tabular}{@{}lll@{}}
\toprule
Concept & Rule & Gotcha\\
\midrule
Sharpe & $\mathrm{se}\sim\sqrt{(1+\mathrm{SR}^2/2)/T}$ & 1yr $\Rightarrow$ CI $\pm1$\\
MDD & worst peak-to-trough & path-aware, strategy-dependent\\
Skew/worst & tail shape & high SR can hide a cliff\\
Hit rate & payoff frequency & meaningless alone\\
Attribution & signal+timing+sizing+cost & net $\neq$ marketable\\
\bottomrule
\end{tabular}
\end{table}

\begin{example}[One monster year inside a flat tape]\label{ex:yrsharpe}
\textbf{Strategy.} Split the real BTC daily series by calendar
year and compute Sharpe, return and MDD per slice. Then compare
with the full-window number you would have quoted.

\textbf{Solution.} \texttt{yearly\_sharpe.py} prints:
\begin{verbatim}
2024: n=71   sharpe +4.08  ret  +36.3%  mdd -13.3%
2025: n=365  sharpe -0.20  ret  -11.8%  mdd -33.5%
2026: n=284  sharpe -0.20  ret  -10.0%  mdd -41.1%
full window: sharpe +0.29 -- the average of very different years
\end{verbatim}
The whole sample's $+0.29$ is an average of a $+4.08$ regime and
two flat $-0.20$ regimes --- the number exists in neither.

\textbf{Significance.} A single Sharpe summarises a world that
never happened. Per-year (or per-regime) slices are the minimal
honesty: if your edge lives in one slice, your backtest is a
study of that slice, not of the market.
\end{example}

\section{Common mistakes}
\begin{itemize}
\item \textbf{``Sharpe 1.8 on the year.''} The CI spans roughly $\pm1$; the claim is
  compatible with mediocre.
\item \textbf{MDD reported without dates.} A $-30\%$ dd lasting 3 years is a different
  product than one lasting 3 weeks.
\item \textbf{Ignoring skew.} Symmetric-vol metrics bless put-sellers.
\item \textbf{Confusing hit rate with edge.} Win frequency and expectancy are different
  axes.
\item \textbf{Skipping turnover.} Return per unit of trading is the honest efficiency.
\item \textbf{Announcing the gross, shipping the net.} Costs and sizing are part of the
  strategy.
\item \textbf{Reporting the pooled Sharpe only.}
$+0.29$ on the full window is the average of $+4.08$ and
$-0.20$ (Example~\ref{ex:yrsharpe}) --- a summary of a regime
mixture, not of the strategy.
\item \textbf{Forgetting the units on MDD.} ``$-33\%$'' needs a
window and a compounding convention --- intraday marks and close
marks give different numbers on the same tape.\end{itemize}

\section{Exercises}
\begin{exercise}[easy]
What is the approximate 95\% CI on measured Sharpe at $T{=}252$ days?
\end{exercise}
\begin{solution}
Theory $\pm0.15$ at SR${=}1$; empirically closer to $\pm1$ on real (non-normal) returns.
\end{solution}

\begin{exercise}[easy]
Why can two strategies share a hit rate and differ everywhere else?
\end{exercise}
\begin{solution}
Hit rate counts sign, not size---it ignores both tail shape and path.
\end{solution}

\begin{exercise}[easy]
Name the metric that catches a put-selling profile.
\end{exercise}
\begin{solution}
Skew (and worst-day): the rare large loss is invisible to symmetric volatility.
\end{solution}

\begin{exercise}[easy]
What are the four attribution categories?
\end{exercise}
\begin{solution}
Signal, timing, sizing, costs.
\end{solution}

\begin{exercise}[easy]
Why is turnover an evaluation metric, not just an execution detail?
\end{exercise}
\begin{solution}
It converts directly to cost and capacity---a 5bp edge trading daily is a different
product than one trading monthly.
\end{solution}

\begin{exercise}
Compute $\mathrm{se}(\mathrm{SR})$ for SR${=}1.5$, $T{=}500$.
\end{exercise}
\begin{solution}
$\sqrt{(1+1.125)/500}\approx0.065$---still a CI a third of a Sharpe wide.
\end{solution}

\begin{exercise}
A strategy: Sharpe 2.5, skew $-5$, MDD $-40\%$. Sell or bury it?
\end{exercise}
\begin{solution}
Bury: the Sharpe is built on a staircase-and-trapdoor profile; the risk model should
price the cliff, not the average.
\end{solution}

\begin{exercise}
How does MDD interact with the evaluation's time horizon?
\end{exercise}
\begin{solution}
Longer horizons mechanically produce deeper MDDs---compare like-for-like windows.
\end{solution}

\begin{exercise}
Write the attribution identity for a backtest return.
\end{exercise}
\begin{solution}
$R_{\text{net}}=R_{\text{signal}}+R_{\text{timing}}+R_{\text{sizing}}-C_{\text{costs}}$.
\end{solution}

\begin{exercise}
Why does skew of $-4$ not show in the Sharpe?
\end{exercise}
\begin{solution}
Volatility is symmetric: $\sigma$ weights up and down moves identically; the cliff is a
third-moment property.
\end{solution}

\begin{exercise}
A backtest: +40\% signal, $-10\%$ sizing, $-8\%$ costs. What do you fix first?
\end{exercise}
\begin{solution}
Sizing---it is implementation, not alpha; costs are a capacity ceiling for later.
\end{solution}

\begin{exercise}
When is hit rate genuinely informative?
\end{exercise}
\begin{solution}
Together with average win/loss ratio---it is half the expectancy identity, useless alone.
\end{solution}

\begin{exercise}[hard]
Show that even at $T{=}2500$ (10y) you cannot separate SR${=}1.0$ from SR${=}0.8$ at
95\% confidence.
\end{exercise}
\begin{solution}
CI $\approx\pm0.05$: the two means differ by 0.2, but each estimate's band is
$[0.95,1.05]$ vs $[0.75,0.85]$---they just separate; in practice fatter tails widen both.
\end{solution}

\begin{exercise}[hard]
Design a metric battery for a put-selling fund---which three numbers dominate and why?
\end{exercise}
\begin{solution}
Worst-day, skew, and recovery time: all three describe the cliff the Sharpe pretends is
volatility.
\end{solution}

\begin{exercise}[hard]
A strategy's entire 5-year MDD occurred in its first year; argue whether the metric is
informative or misleading, and what you'd compute instead.
\end{exercise}
\begin{solution}
Misleading as a forward estimate if the regime changed; compute rolling max-dd or
regime-conditioned dd to see the current risk, not the historical scar.
\end{solution}

\begin{exercise}[easy]
The full-window Sharpe is $+0.29$; two of three slices are
$-0.20$. Reconcile in one line.
\end{exercise}
\begin{solution}
The pooled Sharpe is a weighted average across regimes --- one
short, strong year pulls it up while contributing little to the
per-slice estimates.
\end{solution}
\begin{exercise}[medium]
2024 shows $+4.08$ on $n=71$. Quote an approximate 95\%
confidence interval for that Sharpe (se $\approx1/\sqrt{n}$).
\end{exercise}
\begin{solution}
$4.08\pm1.96/\sqrt{71}\approx4.08\pm0.23$ ---
$[+3.85,+4.31]$... but the regime itself is the uncertainty: a
71-day window cannot certify a long-run property.
\end{solution}
\begin{exercise}[hard]
Design the table that reports these numbers honestly: which
columns and which row tells the truth faster?
\end{exercise}
\begin{solution}
Per-slice rows (n, sharpe, ret, mdd) plus a pooled row marked
``mixture, not strategy'' --- the slice sd (or min--max range)
is the honest summary, the pooled row the cautionary one.
\end{solution}

\section*{Notes}
Evaluation is where optimism goes to be measured. The suite is cheap to compute and
expensive to skip: a Sharpe with no CI, no drawdown, no skew, and no attribution is a
fund's marketing deck, not a research result.

\section*{Further reading}
\begin{itemize}
\item A.~Lo, ``The Statistics of Sharpe Ratios'' --- sampling error, the paper:

\url{https://alo.mit.edu/wp-content/uploads/2017/06/The-Statistics-of-Sharpe-Ratios.pdf}
\item E.~Chan, \emph{Quantitative Trading} --- performance metrics in practice:

\url{https://epchan.blogspot.com/}
\item L.~de Prado lectures --- evaluation metrics and DSR:

\url{https://www.quantresearch.org/Lectures.htm}
\item Bailey \& L\'opez de Prado, ``The Deflated Sharpe Ratio'':

\url{https://www.davidhbailey.com/dhbpapers/deflated-sharpe.pdf}
\item Grinold \& Kahn, \emph{Active Portfolio Management} --- the attribution frame:

\url{https://www.amazon.com/Active-Portfolio-Management-Quantitative-Controlling/dp/0070248826}
\end{itemize}

\begin{rulebox}[Watch instead]
\begin{itemize}
\item ``12. Risk Analysis'' --- MIT 18.S096 (80:32). Risk measures; re-verified 2026-10-11.
\url{https://www.youtube.com/watch?v=qlytPllimpQ}
\item ``11. Time Series Analysis II'' --- MIT 18.S096 (83:47). Testing; re-verified
2026-10-11.
\url{https://www.youtube.com/watch?v=_IFUfFuyQlU}
\end{itemize}
\end{rulebox}
